\documentclass{article}
\usepackage{amsmath,amssymb,amsthm,mathtools}
\newtheorem{theorem}{Theorem}
\newtheorem{lemma}{Lemma}
\newtheorem{proposition}{Proposition}
\newtheorem{corollary}{Corollary}
\newtheorem{definition}{Definition}
\newtheorem{example}{Example}
\title{ShadowBench analysis/L2/ana\_gen\_L2\_002}
\begin{document}
\maketitle
% Problem id: analysis/L2/ana_gen_L2_002
% Companion instructions: docs/instructions.md
% Candidate skeletons: docs/skeletons/
% Lean target: ShadowBench/Source/Main.lean

Theorem (Young's inequility)

If $a,b \ge 0$ are non-negative real numbers and $p,q$ are positive real numbers satisfying $\frac{1}{p} + \frac{1}{q} = 1$, then
$$ ab \le \frac{a^p}{p} + \frac{b^q}{q}$$

proof

Since $p^{-1} + q^{-1} = 1$, by the weighted AM-GM inequality, we have
$$ {a^p}^{p^{-1}}{b^q}^{q^{-1}} \le p^{-1}a^p + q^{-1} b^q$$
which is equivalent to
$$ ab \le \frac{a^p}{p} + \frac{b^q}{q}$$
\end{document}
